\documentclass[11pt]{scrartcl}

\usepackage[top=1.5cm]{geometry}
\usepackage{float}

\setlength{\parindent}{0em}
\setlength{\parskip}{0.5em}

\newcommand{\youranswerhere}{[Your answer goes here \ldots]}
\renewcommand{\thesubsection}{\arabic{subsection}}

\title{
  \textbf{\large Assignment 5} \\
  Join Tuning \\
  {\large Database Tuning}}

\author{
  Group Name (e.g. A1, B5, B3) \\
  \large Lastname1 Firstname1, StudentID1 \\
  \large Lastname2 Firstname2, StudentID2 \\
  \large Lastname3 Firstname3, StudentID3
}

\begin{document}

\maketitle

\subsection*{Experimental Setup}

Describe your experimental setup in a few lines.

\youranswerhere{}

\subsection*{Join Strategies Proposed by System}

\paragraph{Response times}\mbox{}

\begin{table}[H]
  \centering
  \begin{tabular}{l|l|l}
    Indexes & Join Strategy Q1 & Join Strategy Q2 \tabularnewline
    \hline
    no index & 1716.181 & 422.025 \tabularnewline
    unique non-clustering on \texttt{Publ.pubID} & 1681.725 & 112.743
      \tabularnewline
    clustering on \texttt{Publ.pubID} and \texttt{Auth.pubID} & 1082.134 & 257.781
      \tabularnewline
  \end{tabular}
\end{table}

\paragraph{Discussion}
\textbf{No index.}
For Q1 the system chooses a \emph{Hash Join}: it scans \texttt{Publication} sequentially, builds a hash table on \texttt{pubID}, and then probes it while scanning \texttt{Author} sequentially.
This is the expected choice when no index is available and the join produces millions of output rows, because a hash join runs in $O(n+m)$ and only requires one pass over each relation (assuming the hash table fits in memory).
For Q2 the system uses a \emph{Parallel Hash Join}: the build side is the filtered \texttt{Author} scan (very few rows after the \texttt{name} filter), and the probe side is a parallel sequential scan of \texttt{Publication}.
Because the filter on \texttt{Auth.name} is applied before the join and reduces the inner side dramatically, parallelism further reduces wall-clock time.
Without any index the planner cannot avoid reading all author pages to evaluate the name predicate, so a sequential scan with a filter remains the only option.

\textbf{Unique non-clustering index on \texttt{Publ.pubID}.}
For Q1 the plan is again a \emph{Hash Join} identical to the no-index case.
Because Q1 retrieves all publications (no selectivity on the \texttt{Publication} side), using the non-clustering index would require one random I/O per row — far more expensive than a single sequential scan — so the planner correctly ignores the index.
For Q2 the system switches to a \emph{Nested Loop}: it first scans \texttt{Author} sequentially with the \texttt{name} filter and then, for each of those tuples, performs an \emph{Index Scan} on \texttt{Publication} via \texttt{temp\_nonc\_publ} to retrieve the corresponding publication.
This is the expected choice: with a highly selective outer relation the nested loop avoids building any hash table, and each index lookup costs only $O(\log n)$.
The non-clustering index is sufficient here because only a small number of random page accesses are needed (one per author tuple of `Divesh Srivastava').

\textbf{Clustering indexes on \texttt{Publ.pubID} and \texttt{Auth.pubID}.}
For Q1 the system selects a \emph{Merge Join}: both tables are accessed via their respective clustering index scans (\texttt{temp\_c\_publ} and \texttt{temp\_c\_auth}), which deliver rows already sorted on \texttt{pubID}.
A merge join over two pre-sorted inputs runs in $O(n+m)$ without any additional sort step, and the clustering indexes ensure sequential I/O on both sides, making this the most efficient strategy for the large Cartesian-like join.
For Q2 the planner again chooses a \emph{Nested Loop} with a parallel sequential scan of \texttt{Author} (filter on \texttt{name}) as the outer relation and an \emph{Index Scan} on \texttt{Publication} as the inner.
Even though a clustering index on \texttt{Auth.pubID} exists, scanning \texttt{Author} sequentially with a filter is cheaper than using that index to find `Divesh Srivastava' entries (the name column is not indexed).
Once the few matching author tuples are found, the clustering index on \texttt{Publ.pubID} enables efficient point lookups, keeping the total cost low.
The observed response times confirm these choices: Q1 benefits most from the merge join (1082\,ms vs.\ $\approx$1700\,ms without clustering), while Q2 is fastest with the non-clustering index variant (113\,ms) because the index lookup cost is the same but the non-clustering index case avoids the overhead of the parallel gather node present in the clustering variant.

No Index:

{\small
\parskip0pt\begin{verbatim}
[
Hash Join (cost=68240.32..235012.34 rows=3095201 width=82)
Hash Cond: ((a.pubid)::text = (p.pubid)::text)
-> Seq Scan on author a (cost=0.00..57784.01 rows=3095201 width=38)
-> Hash (cost=34760.14..34760.14 rows=1233214 width=89)
-> Seq Scan on publication p (cost=0.00..34760.14 rows=1233214 width=89)

Gather (cost=43952.96..73448.69 rows=24 width=82)
Workers Planned: 2
-> Parallel Hash Join (cost=42952.96..72446.29 rows=10 width=82)
Hash Cond: ((p.pubid)::text = (a.pubid)::text)
-> Parallel Seq Scan on publication p (cost=0.00..27566.39 rows=513839 width=89)
-> Parallel Hash (cost=42952.84..42952.84 rows=10 width=38)
-> Parallel Seq Scan on author a (cost=0.00..42952.84 rows=10 width=38)
Filter: ((name)::text = 'Divesh Srivastava'::text)
]
\end{verbatim}}

Unique Non-clustering index on Publ.pubID:

{\small
\parskip0pt\begin{verbatim}
[
Hash Join (cost=68240.32..200578.26 rows=3095201 width=82)
Hash Cond: ((a.pubid)::text = (p.pubid)::text)
-> Seq Scan on author a (cost=0.00..57784.01 rows=3095201 width=38)
-> Hash (cost=34760.14..34760.14 rows=1233214 width=89)
-> Seq Scan on publication p (cost=0.00..34760.14 rows=1233214 width=89)

Gather (cost=1000.43..44039.69 rows=24 width=82)
Workers Planned: 2
-> Nested Loop (cost=0.43..43037.29 rows=10 width=82)
-> Parallel Seq Scan on author a (cost=0.00..42952.84 rows=10 width=38)
Filter: ((name)::text = 'Divesh Srivastava'::text)
-> Index Scan using temp_nonc_publ on publication p (cost=0.43..8.45 rows=1 width=89)
Index Cond: ((pubid)::text = (a.pubid)::text)
]
\end{verbatim}}

Clustering index on Publ.pubID and Auth.pubID:

{\small
\parskip0pt\begin{verbatim}
[
Merge Join (cost=0.86..218772.58 rows=3095201 width=82)
Merge Cond: ((p.pubid)::text = (a.pubid)::text)
-> Index Scan using temp_c_publ on publication p (cost=0.43..67009.64 rows=1233214 width=89)
-> Index Scan using temp_c_auth on author a (cost=0.43..110152.45 rows=3095201 width=38)

Gather (cost=1000.43..44016.79 rows=24 width=83)
Workers Planned: 2
-> Nested Loop (cost=0.43..43014.39 rows=10 width=83)
-> Parallel Seq Scan on author a (cost=0.00..42929.84 rows=10 width=38)
Filter: ((name)::text = 'Divesh Srivastava'::text)
-> Index Scan using temp_c_publ on publication p (cost=0.43..8.45 rows=1 width=90)
Index Cond: ((pubid)::text = (a.pubid)::text)
]
\end{verbatim}}

\subsection*{Indexed Nested Loop Join}

\paragraph{Response times}\mbox{}

\begin{table}[H]
  \centering
  \begin{tabular}{l|r|r}
    Indexes & Response time Q1 [ms] & Response time Q2 [ms] \tabularnewline
    \hline
    index on \texttt{Publ.pubID} & \ldots & \ldots \tabularnewline
    index on \texttt{Auth.pubID} & \ldots & \ldots \tabularnewline
    index on \texttt{Publ.pubID} and \texttt{Auth.pubID} & \ldots & \ldots
        \tabularnewline
  \end{tabular}
\end{table}

\paragraph{Query plans}\mbox{}

Index on \texttt{Publ.pubID} (Q1/Q2):

{\small
\parskip0pt\begin{verbatim}
[Your query plans (index on Publ.pubID) go here ...]
\end{verbatim}}

Index on \texttt{Auth.pubID} (Q1/Q2):

{\small
\parskip0pt\begin{verbatim}
[Your query plans (index on Auth.pubID) go here ...]
\end{verbatim}}

Index on \texttt{Auth.pubID} and \texttt{Auth.pubID} (Q1/Q2):

{\small
\parskip0pt\begin{verbatim}
[Your query plans (index on Publ.pubID and Auth.pubID) go here ...]
\end{verbatim}}

\paragraph{Discussion}

Discuss your observations. Are the response times expected? Why (not)?

\youranswerhere{}

\subsection*{Sort-Merge Join}

\paragraph{Response times}\mbox{}

\begin{table}[H]
  \centering
  \begin{tabular}{l|r|r}
    Indexes & Response time Q1 [ms] & Response time Q2 [ms] \tabularnewline
    \hline
    no index & \ldots & \ldots \tabularnewline
    two non-clustering indexes & \ldots & \ldots \tabularnewline
    two clustering indexes & \ldots & \ldots \tabularnewline
  \end{tabular}
\end{table}

\paragraph{Query plans}\mbox{}

No index (Q1/Q2):

{\small
\parskip0pt\begin{verbatim}
[Your query plans (no index) go here ...]
\end{verbatim}}

Two non-clustering indexes (Q1/Q2):
{\small
\parskip0pt\begin{verbatim}
[Your query plans (two non-clustering indexes) go here ...]
\end{verbatim}}

Two clustering indexes  (Q1/Q2):
{\small
\parskip0pt\begin{verbatim}
[Your query plans (two clustering indexes) go here ...]
\end{verbatim}}

\paragraph{Discussion}

Discuss your observations. Are the response times expected? Why (not)?

\youranswerhere{}

\subsection*{Hash Join}

\paragraph{Response times}\mbox{}

\begin{table}[H]
  \centering
  \begin{tabular}{l|r|r}
    Indexes & Response time Q1 [ms] & Response time [ms] Q2 \tabularnewline
    \hline
    no index & 2647.080 & 101.186 \tabularnewline
  \end{tabular}
\end{table}

\paragraph{Query plans}\mbox{}

No Index (Q1/Q2):

{\small
\parskip0pt\begin{verbatim}
[
Hash Join (cost=68240.32..235012.34 rows=3095201 width=82)
Hash Cond: ((a.pubid)::text = (p.pubid)::text)
-> Seq Scan on author a (cost=0.00..57784.01 rows=3095201 width=38)
-> Hash (cost=34760.14..34760.14 rows=1233214 width=89)
-> Seq Scan on publication p (cost=0.00..34760.14 rows=1233214 width=89)

Gather (cost=43952.96..73448.69 rows=24 width=82)
Workers Planned: 2
-> Parallel Hash Join (cost=42952.96..72446.29 rows=10 width=82)
Hash Cond: ((p.pubid)::text = (a.pubid)::text)
-> Parallel Seq Scan on publication p (cost=0.00..27566.39 rows=513839 width=89)
-> Parallel Hash (cost=42952.84..42952.84 rows=10 width=38)
-> Parallel Seq Scan on author a (cost=0.00..42952.84 rows=10 width=38)
Filter: ((name)::text = 'Divesh Srivastava'::text)
]
\end{verbatim}}

\paragraph{Discussion}

<<<<<<< HEAD
What do you think about the response time of the hash index vs.\ the response times of sort-merge and index nested loop join for each of the queries? Explain.

\youranswerhere{}
=======
The forced hash join was measured with no index present, yielding \textbf{2647\,ms} for Q1 and \textbf{101\,ms} for Q2.

\textbf{Q1 (full join, no selectivity on \texttt{name}).}
The hash join (2647\,ms) is the slowest strategy for this query among all variants tested.
The best sort-merge variant (two non-clustering indexes, no explicit sort needed) completes in 962\,ms, and the best indexed nested loop variant (index on \texttt{Auth.pubID} only) takes 1073\,ms.
Even without any index, sort-merge finishes in 2215\,ms — still faster than the forced hash join.
The reason is that Q1 produces over 3 million output rows, causing the hash table for \texttt{Publication} (1.2M tuples) to spill to disk in 32 batches (as reported by \texttt{Batches: 32}), which generates 16135 temporary pages of I/O.
Sort-merge with indexes avoids this entirely by reading both relations in order, while INLJ with an index on \texttt{Auth.pubID} benefits from probing the smaller publication table from the outer loop.
In short, the hash join's disk-spill overhead makes it the least attractive choice when the join result is very large.

\textbf{Q2 (selective filter on \texttt{Auth.name = 'Divesh Srivastava'}).}
Here the hash join (101\,ms) is competitive and substantially faster than sort-merge under all index configurations (301--409\,ms).
It is also faster than the INLJ variant with only an index on \texttt{Auth.pubID} (926\,ms), although it is slightly slower than INLJ with an index on \texttt{Publ.pubID} or both (55--60\,ms).
The Parallel Hash Join used for Q2 benefits from parallelism (2 workers): the build phase operates only on the tiny filtered author set ($\approx$183 rows), keeping the hash table entirely in memory (no batching), and the probe side scans \texttt{Publication} in parallel.
Sort-merge is slower for Q2 because it must sort (or index-scan) the full \texttt{Publication} relation regardless of the name filter, adding significant I/O even when an index on \texttt{pubID} exists.
The INLJ variant with a \texttt{Publ.pubID} index remains the fastest for Q2 because it probes \texttt{Publication} only for the 183 qualifying author tuples, but the hash join is the better general-purpose choice when no index on \texttt{Publ.pubID} is available.

\subsection*{Time Spent on this Assignment}

Time in hours per person: \textbf{XXX}

\subsection*{References}

\begin{table}[H]
  \centering
  \begin{tabular}{c}
    \hline
    \textbf{Important:} Reference your information sources! \tabularnewline
    Remove this section if you use footnotes to reference your information
    sources. \tabularnewline
    \hline
  \end{tabular}
\end{table}

\end{document}